\documentclass[11pt,a4paper]{article}
\usepackage[T1]{fontenc}
\usepackage[utf8]{inputenc}
\usepackage{lmodern,microtype}
\usepackage[margin=27mm,headheight=15pt]{geometry}
\usepackage{amsmath,amssymb,amsthm,mathtools,bm}
\usepackage{booktabs,array,longtable,enumitem}
\usepackage[dvipsnames]{xcolor}
\usepackage{fancyhdr}
\usepackage{hyperref}
\usepackage[nameinlink,noabbrev]{cleveref}
\hypersetup{colorlinks=true,linkcolor=MidnightBlue,citecolor=MidnightBlue,urlcolor=MidnightBlue,
 pdftitle={Slowly Varying Critical Transseries},pdfauthor={Research article prepared for Vladimir Reshetnikov}}
\pagestyle{fancy}\fancyhf{}
\fancyhead[L]{\small Slowly varying critical transseries}
\fancyhead[R]{\small ProveIt research continuation}
\fancyfoot[C]{\thepage}
\setlength{\parskip}{4pt}\setlength{\parindent}{1.1em}
\setlength{\emergencystretch}{2em}
\numberwithin{equation}{section}
\newtheorem{theorem}{Theorem}[section]
\newtheorem{proposition}[theorem]{Proposition}
\newtheorem{lemma}[theorem]{Lemma}
\newtheorem{corollary}[theorem]{Corollary}
\theoremstyle{definition}\newtheorem{definition}[theorem]{Definition}
\newtheorem{example}[theorem]{Example}
\theoremstyle{remark}\newtheorem{remark}[theorem]{Remark}
% ed. (2026-09-29): unnumbered environment for ProveIt editorial notes; it does not shift any numbering.
\theoremstyle{definition}
\newtheorem*{ednoteinner}{Editorial note (ProveIt, 2026-09-29)}
\newenvironment{ednote}{\begin{ednoteinner}\sloppy\Urlmuskip=0mu plus 1mu\relax}{\end{ednoteinner}}
\DeclareMathOperator{\Li}{Li}
\DeclareMathOperator{\supp}{supp}
\DeclareMathOperator{\Rea}{Re}
\newcommand{\R}{\mathbb R}\newcommand{\C}{\mathbb C}
\newcommand{\N}{\mathbb N}\newcommand{\E}{\mathbb E}
\newcommand{\Prob}{\mathbb P}\newcommand{\ii}{\mathrm i}
\newcommand{\bc}{\beta_{\mathrm c}}\newcommand{\rhoc}{\rho_{\mathrm c}}
\newcommand{\G}{\Gamma(-\alpha)}
\newcommand{\dd}{\,\mathrm d}
\newcommand{\F}{\mathcal F^{-1}}
\newcommand{\asyminv}{\leftarrow}
\newcommand{\repo}{\url{https://github.com/VladimirReshetnikov/ProveIt}}
\title{\vspace{-12mm}\textbf{Slowly Varying Critical Transseries}\\[4mm]
\Large Universal action budgets, convergent Lambert charts,\\
and an effective-index expansion}
\author{Research article prepared for Vladimir Reshetnikov}
\date{September 29, 2026}
\begin{document}
\maketitle
\begin{abstract}
Consider the countable exponential-feedback equation
$U=\beta A(qe^U)$ with nonnegative coefficients
$a_j\sim a j^{-1-\alpha}\ell(j)$, where $1<\alpha<2$ and $\ell$ is slowly
varying. The ProveIt critical-Hahn manuscript explicitly asks whether its
exact-power-tail theory extends to this class. We give a fixed-interior-exponent
answer. The critical inverse is described by an asymptotic inverse of
$y^\alpha/\ell(y)$; the same normalization controls the large coefficients
and the sharp scale of asymptotically sufficient action cutoffs. A compensated
stable density supplies a universal coupling--cutoff profile, independent
of $\ell$ after normalization, and a conditional point-process formula
identifies the entire cloud of large actions.

The representation question has a genuine obstruction: an explicit positive,
slowly varying oscillatory tail yields no leading monomial in any finite
iterated power--log scale. For the exact tails
$j^{-1-\alpha}(\log j)^m$, $m\ge1$, this obstruction disappears. We construct
an absolutely convergent three-coordinate analytic chart over a
$W_{-1}$ core, with explicit first coefficients and a geometric tail bound.
We also prove all-order inverse-logarithmic expansions for coefficients
and cutoff probabilities. Their first correction is exactly the derivative
obtained by replacing $\alpha$ by $\alpha-m/\log b_n$ while holding the
normalized coordinates fixed. A positive derivative formula for the cutoff
profile turns this into a two-term minimal-budget law.

The stable-limit and regular-variation mechanisms are classical. This article supplies explicit inverse charts, differential correction
operators, and action-budget formulas, with proofs connecting them to that
classical theory. No
claim of historical priority, endpoint uniformity, or Lean formalization
is made. The package contains 552 exact algebra checks and separately
labelled floating-point diagnostics.
\end{abstract}
\vspace{2mm}
\noindent\textbf{Keywords.} Transseries; Lagrange inversion; regular variation;
stable local limits; action truncation; Lambert $W$; logarithmic asymptotics.
\clearpage
{\small\setlength{\parskip}{0pt}\tableofcontents}
\newpage

\section{The question being answered and the limits of the answer}
\label{sec:position}

\subsection{Repository provenance}
The repository was inspected at commit
\begin{center}\small\texttt{ebd8344bca77d8352cd745b7df374618a290a029}.\end{center}
The primary predecessor is \emph{Beyond Finite-Action Folds: Critical Hahn
Transseries, Stable Sector Asymptotics, and Sharp Action Budgets}, located
under the transseries documentation directory in
\texttt{Critical\_Hahn\_Transseries\_Beyond\_Finite\_Action\_Folds/}
\cite{repo-critical}. Its research Question~1, ``Slowly varying action tails,''
asks for the implicit stable scale, the cutoff profile, an inverse notation
when de Bruijn conjugates are necessary, and hypotheses preserving a finite
positive support grid. That question is the direct target of this article.

The predecessor already treats exact power tails, a bounded coupling window,
and the conditional-Poisson interpretation of action truncation. These are
not claimed as new here. The repository index also records upper-endpoint
articles involving cubic tails and logarithmic modifications
\cite{repo-index}. Our exponent stays in $(1,2)$; neither an $\alpha\to2$
limit nor the Gaussian endpoint is being repeated or asserted. The repository
index warns that its recent research arrivals are unmerged and not formally
verified. We therefore prove the mathematical statements needed below rather
than use those arrivals as unexamined analytic axioms.
% ed. (2026-09-29): endpoint packages in (or filed after) this article's snapshot that it does not cite by name.
\begin{ednote}
The logarithmic-endpoint package
(\path{Logarithmic_Critical_Endpoint_Lambert_Charts/} \cite{ed:lce}, in the
tree at the inspected snapshot) treats the $\alpha=2$ analogue of the tails
of \cref{sec:chart}, $a_j\sim aj^{-3}(\log j)^r$, $r\ge-1$, at leading
order (its \texttt{thm:weightedchart}, \texttt{thm:weightedbudget}); its
Lambert core has the form of \eqref{eq:Lambertcore} at $(\alpha,m)=(2,1)$,
where the logarithm comes from the pole of $\Gamma(-\alpha)$. The marginal
package (\path{Marginal_Critical_Transseries_Action_Budgets/} \cite{ed:mct},
\texttt{thm:min-budget}) gives the $\alpha=2$ two-term minimal budget; the
later confluent and stable--Gaussian packages \cite{ed:cct,ed:sge} treat the
crossover $\alpha\to2$ for exact tails, where the cutoff profile tends to
$\exp(-s^{-\alpha}/\alpha)$. None of them treats $1<\alpha<2$ with a slowly
varying factor.
\end{ednote}

\subsection{Results and their logical scope}
For arbitrary positive slowly varying $\ell$, we prove a critical real inverse
law, a stable local limit, a two-parameter cutoff profile, and the exact
criterion $M_n/b_n\to\infty$ for asymptotically lossless truncation in a bounded
critical window. These conclusions require no complex continuation of $\ell$.
An oscillatory example shows why they cannot automatically be upgraded to a
power--logarithmic transseries representation.

For the smaller, analytically controlled class with eventual exact weights
$a j^{-1-\alpha}(\log j)^m$, we prove much more: a convergent local inverse
chart, all fixed orders of a coefficient expansion, the corresponding
cutoff expansion, and an asymptotic minimal-budget formula. The coefficient
expansion is asymptotic, not asserted to converge; the local inverse chart
really does converge. Confusing those two statements would erase the main
analytic distinction of the paper.

\subsection{Classical input and novelty discipline}
Regular variation, asymptotic inversion, and de Bruijn conjugacy are part of
the classical theory developed in \cite{bgt}. The relation between generating
functions, conditioned allocations, stable limits, and largest components
is classical as well; Janson's survey is a detailed reference
\cite{janson}. Our probability proof is a direct Fourier proof for this
specific compound-Poisson model, not a claim to have invented a stable local
limit theorem. The convergent special-function expansion used for exact
logarithmic tails is the parameter-differentiated polylogarithm expansion
in \cite{dlmf-polylog}; the Lambert branch is standard \cite{dlmf-lambert}.
Transseries terminology and the difference between formal operations and
analytic realization are discussed in \cite{edgar}.

The principal synthesis developed here is the following: the same slowly
varying normalization controls inverse singularities and action budgets;
logarithmic tails give an explicit convergent inverse atlas; and logarithmic
perturbations of the L\'evy measure induce a common differential correction
on the density, the retained fraction, and its quantiles. These are results
proved in this article, not independently certified claims of publication
priority. The literature comparison is targeted rather than exhaustive.

\section{Model, normalization, and exact coefficient identities}
\label{sec:model}

Fix $1<\alpha<2$, $a>0$, and a sequence $(a_j)_{j\ge1}$ satisfying
\begin{equation}\label{eq:assumptions}
 a_j\ge0,\qquad a_1>0,\qquad
 a_j\sim a j^{-1-\alpha}\ell(j),
\end{equation}
where $\ell$ is eventually positive, measurable, locally bounded, and slowly
varying at infinity:
\[
 \frac{\ell(cy)}{\ell(y)}\longrightarrow1
 \quad\text{for every fixed }c>0.
\]
Values of $\ell$ on an initial bounded interval are immaterial. Set
\begin{equation}\label{eq:constants}
 A(t)=\sum_{j\ge1}a_jt^j,\quad d_0=A(1),\quad
 d_1=\sum_{j\ge1}j a_j,\quad
 \bc=d_1^{-1},\quad \rho_\beta=e^{-\beta d_0}.
\end{equation}
Both $d_0$ and $d_1$ are finite and positive. The second moment is infinite.
For $\beta>0$ define
\begin{equation}\label{eq:feedback}
 U_\beta(q)=\beta A(qe^{U_\beta(q)})
           =\sum_{n\ge1}u_n(\beta)q^n.
\end{equation}
A formal coefficient recursion and the analytic implicit-function theorem
at $(q,U)=(0,0)$ give the same unique analytic germ. Nonnegativity and $a_1>0$
give $u_n(\beta)>0$ for all $n\ge1$.

For a positive integer $M$ replace $A$ by
$A_M(t)=\sum_{j\le M}a_jt^j$ and write $u_{n,M}(\beta)$ for the resulting
coefficients. This is a literal removal of actions, without renormalizing
the remaining weights or changing $\beta$.

\begin{proposition}[Exact coefficient and conditioning identities]
\label{prop:exact}
For $n\ge1$,
\begin{equation}\label{eq:lagrange}
 u_n(\beta)=\frac1n[t^n]e^{n\beta A(t)}.
\end{equation}
Let $N_{n,j}$ be independent Poisson variables of means $n\beta a_j$ and
$K_n=\sum_{j\ge1}jN_{n,j}$. Then
\begin{align}
 u_n(\beta)&=\frac{\rho_\beta^{-n}}n\Prob(K_n=n),\label{eq:poisson}\\
 \frac{u_{n,M}(\beta)}{u_n(\beta)}
 &=\Prob(N_{n,j}=0\text{ for every }j>M\mid K_n=n).\label{eq:ratio}
\end{align}
The ratio is nondecreasing in $M$, belongs to $[0,1]$, and equals $1$ for
$M\ge n$.
\end{proposition}
\begin{proof}
Put $t=qe^{U_\beta}$, so $t=q e^{\beta A(t)}$ and $U_\beta=\beta A(t)$.
Lagrange inversion gives
\[
 [q^n]U_\beta=\frac1n[t^{n-1}]\beta A'(t)e^{n\beta A(t)}
             =\frac1n[t^n]e^{n\beta A(t)}.
\]
The total Poisson intensity is $n\beta d_0<\infty$, so only finitely many
$N_{n,j}$ are nonzero almost surely. Its generating function is
$\E t^{K_n}=\exp(n\beta(A(t)-d_0))$, proving \eqref{eq:poisson}.
On the no-large-action event, independence leaves
$\exp(n\beta(A_M(t)-d_0))$. Comparing its coefficient with the full one
proves \eqref{eq:ratio}. Its remaining properties follow from nested events.
\end{proof}

\begin{definition}[Critical scale and coupling coordinate]
\label{def:scale}
Let $\beta_n\to\bc$ and choose $b_n\to\infty$ such that
\begin{equation}\label{eq:bscale}
 n\beta_n a\,\ell(b_n)b_n^{-\alpha}\longrightarrow1.
\end{equation}
Define
\begin{equation}\label{eq:lambdan}
 \lambda_n=\frac{n(1-\beta_n d_1)}{b_n}.
\end{equation}
Unless stated otherwise, the coupling coordinate is confined to a fixed
compact subset of $\R$.
\end{definition}
Such a scale exists by asymptotic inversion of a regularly varying function
of positive index. It is unique up to asymptotic equivalence and has index
$1/\alpha$ as a function of $n$. In particular $b_n=o(n)$. Replacing it by an
asymptotically equivalent scale changes neither the leading laws nor the
criterion for a lossless cutoff. Higher-order laws will require a more
precise normalization.

\section{The critical inverse for an arbitrary slowly varying tail}
\label{sec:inverse-rv}

Write $t=e^{-x}$ and let $q=t e^{-\bc A(t)}$. The exact real boundary
coordinate and inverse value are
\begin{align}
 \varepsilon&=\log(\rhoc/q)
       =x+\bc\bigl(A(e^{-x})-d_0\bigr),\label{eq:epsilon}\\
 U_{\bc}(q)&=\bc d_0+\varepsilon-x.\label{eq:Uinverse}
\end{align}
Here $x\downarrow0$ as $q\uparrow\rhoc$. The derivative of the right side
of \eqref{eq:epsilon} is
$1-\bc\sum j a_j e^{-jx}>0$ for $x>0$. Thus the positive inverse exists
and is unique near the boundary.

\begin{lemma}[Compensated Abelian limit]\label{lem:abelian}
Under \eqref{eq:assumptions},
\begin{equation}\label{eq:abelian}
 A(e^{-x})-d_0+d_1x
 \sim a\Gamma(-\alpha)x^\alpha\ell(1/x)
 \qquad(x\downarrow0).
\end{equation}
\end{lemma}
\begin{proof}
Put $h(v)=e^{-v}-1+v$. The expression on the left is
$\sum a_jh(jx)$. On $v$ in a compact subinterval of $(0,\infty)$,
regular variation converts its normalization by
$a x^\alpha\ell(1/x)$ into the Riemann sum for
$\int h(v)v^{-1-\alpha}\dd v$. The uniform convergence theorem and Potter
bounds for slowly varying functions justify the passage from compact
subintervals to the full half-line \cite{bgt}. For completeness, choose
$0<\eta<\min(2-\alpha,\alpha-1)$. Near zero use $h(v)\le v^2/2$ and the
bound $C v^{-1-\alpha-\eta}$ for the normalized weight; near infinity use
$h(v)\le v$ and $C v^{-1-\alpha+\eta}$. The two dominating functions are
integrable. An initial finite prefix contributes $O(x^2)$, which is
$o(x^\alpha\ell(1/x))$. Two integrations by parts give
\[
 \int_0^\infty h(v)v^{-1-\alpha}\dd v
 =\frac{\Gamma(2-\alpha)}{\alpha(\alpha-1)}
 =\Gamma(-\alpha)>0.
\]
All boundary terms vanish by $1<\alpha<2$.
\end{proof}

Let $R(y)=y^\alpha/\ell(y)$. We use $R^\asyminv$ to denote any asymptotic
inverse, not necessarily the ordinary inverse of the original possibly
nonmonotone representative. One can first replace $R$ by an asymptotically
equivalent increasing function.

\begin{theorem}[Critical inverse and conjugate notation]\label{thm:rv-inverse}
Put $C=\bc a\Gamma(-\alpha)$. The solution of \eqref{eq:epsilon} satisfies
\begin{equation}\label{eq:inverse-rv}
 \boxed{\quad x(\varepsilon)\sim
 \frac1{R^\asyminv(C/\varepsilon)}.\quad}
\end{equation}
In particular $x$ is regularly varying at zero with index $1/\alpha$.
If $k(y)=\ell(y)^{-1/\alpha}$ and $k^\#$ is a de Bruijn conjugate, then,
with $z=(C/\varepsilon)^{1/\alpha}$,
\begin{equation}\label{eq:conjugate}
 x(\varepsilon)\sim\frac1{z k^\#(z)},\qquad
 b_n\sim(n\beta_n a)^{1/\alpha}
 k^\#\bigl((n\beta_n a)^{1/\alpha}\bigr).
\end{equation}
\end{theorem}
\begin{proof}
By \cref{lem:abelian}, $\varepsilon\sim C/R(1/x)$. Asymptotic inversion
of a regularly varying function with positive index proves
\eqref{eq:inverse-rv}. The conjugacy relations are
\[
 k(y)k^\#(yk(y))\to1,\qquad
 k^\#(y)k(yk^\#(y))\to1.
\]
Since $R(y)=(yk(y))^\alpha$, the second relation implies
$R(zk^\#(z))\sim z^\alpha$. This proves the first formula in
\eqref{eq:conjugate}, and the same argument applied to \eqref{eq:bscale}
proves the second. Existence of a conjugate and uniqueness up to asymptotic
equivalence are the standard conjugacy theorem of \cite{bgt}.
\end{proof}

\begin{corollary}[Reciprocal matching]\label{cor:reciprocal}
At $\beta_n=\bc$,
\[
 \varepsilon(1/b_n)\sim\frac{\Gamma(-\alpha)}n.
\]
Thus the inverse singularity scale and the large-coefficient fluctuation
scale are reciprocal, with the displayed gamma normalization.
\end{corollary}
\begin{proof}
Substitute $x=1/b_n$ into \cref{lem:abelian} and use \eqref{eq:bscale}.
\end{proof}

This theorem is a real asymptotic statement. Slow variation alone does not
supply analytic continuation of $A$ across its unit-circle boundary, a
complex inverse atlas, or a convergent Hahn expansion. The following
counterexample shows that even the leading monomial can fail.

\section{A genuine obstruction to power--logarithmic representation}
\label{sec:obstruction}

\begin{theorem}[A slowly varying tail without a power--log leading monomial]
\label{thm:obstruction}
Take, for all sufficiently large $j$,
\begin{equation}\label{eq:oscillating}
 a_j=a j^{-1-\alpha}\bigl(2+\sin(\log\log j)\bigr),
\end{equation}
and choose any nonnegative finite prefix with $a_1>0$. Its positive critical
inverse $x(\varepsilon)$ is not asymptotic to any expression
\begin{equation}\label{eq:plmonomial}
 c\,\varepsilon^r
 \prod_{h=1}^d\bigl(\log_h(1/\varepsilon)\bigr)^{s_h},
 \qquad c>0,\quad r,s_h\in\R,\quad d<\infty,
\end{equation}
where $\log_h$ denotes the $h$-fold iterated logarithm. Consequently it has
no convergent well-based expansion in these monomials whose first nonzero
term is asymptotically dominant. In particular no ordinary finite Puiseux
chart exists.
\end{theorem}
\begin{proof}
The factor $\ell(y)=2+\sin(\log\log y)$ is positive and slowly varying:
for fixed $c>0$, the change in its sine argument is $O(1/\log y)$, and
$\ell(y)\ge1$. By \cref{thm:rv-inverse}, any monomial equivalent must have
$r=1/\alpha$. The inequalities $1\le\ell\le3$ and
\cref{lem:abelian} imply that $x(\varepsilon)/\varepsilon^{1/\alpha}$
is bounded above and below by positive constants. In a finite product of
iterated logarithmic powers, the first nonzero exponent decides whether
the product tends to zero or infinity. Therefore all $s_h$ in
\eqref{eq:plmonomial} must vanish.

It remains to disprove convergence to a constant. Let
\[
 x_k=\exp\{-\exp(2\pi k+\pi/2)\},\qquad
 y_k=\exp\{-\exp(2\pi k+3\pi/2)\}.
\]
Along these two sequences $\ell(1/x_k)=3$ and $\ell(1/y_k)=1$. With
$C=\bc a\Gamma(-\alpha)$, \cref{lem:abelian} yields
\[
 \frac{x_k}{\varepsilon(x_k)^{1/\alpha}}\to(3C)^{-1/\alpha},\qquad
 \frac{y_k}{\varepsilon(y_k)^{1/\alpha}}\to C^{-1/\alpha}.
\]
The limits differ. A well-based expansion with a dominant leading term
would provide the forbidden monomial equivalent.
\end{proof}

\begin{remark}[What the obstruction does not say]
The theorem concerns the specified finite iterated power--log scales and
expansions with their usual leading-term meaning. It does not exclude a
calculus enlarged by oscillatory slow blocks, and it is not a theorem that
all possible generalized transseries fields fail. It does show that
``slowly varying'' is not a sufficient hypothesis for the repository's
finite positive power--log grid to survive unchanged.
\end{remark}

\section{Stable local limits from a direct Fourier argument}
\label{sec:local}

The normalization throughout is the L\'evy measure
\begin{equation}\label{eq:levy}
 \nu_\alpha(\dd x)=x^{-1-\alpha}\mathbf1_{x>0}\dd x.
\end{equation}
Define
\begin{align}
 J_\alpha(t)
 &=\int_0^\infty(e^{\ii tx}-1-\ii tx)x^{-1-\alpha}\dd x
   =\Gamma(-\alpha)(-\ii t)^\alpha,\label{eq:J}\\
 g_\alpha(y)&=\frac1{2\pi}\int_\R
              e^{-\ii ty+J_\alpha(t)}\dd t.\label{eq:g}
\end{align}
The complex power uses the principal branch. Since
$\cos(\pi\alpha/2)<0$, the Fourier integral converges absolutely and
remains so after multiplication by any polynomial in $t$ or $\log|t|$.
The corresponding probability law can be constructed directly: compensate
the Poisson jumps on $(0,1]$ and pass to their $L^2$ limit, using
$\int_0^1 x^2\nu_\alpha(\dd x)<\infty$; then add the independent,
finite-intensity jumps on $(1,\infty)$ and subtract their finite mean.
Its characteristic exponent is $J_\alpha$. The same construction with an
upper jump cutoff defines the bounded-jump laws used below.
At zero,
\begin{equation}\label{eq:g0}
 g_\alpha(0)=\frac{\Gamma(-\alpha)^{-1/\alpha}}
 {\alpha\Gamma(1-1/\alpha)}
 =\frac{\Gamma(-\alpha)^{-1/\alpha}}{|\Gamma(-1/\alpha)|}.
\end{equation}
One obtains this directly by integrating
$\exp\{-z t^\alpha\}$ on $t>0$ with
$z=-\Gamma(-\alpha)e^{-\ii\pi\alpha/2}$ and applying the reflection formula.

\begin{lemma}[Scaled measure convergence and Fourier damping]
\label{lem:fourier}
Let \eqref{eq:bscale} hold and put
$\nu_n=n\beta_n\sum_{j\ge1}a_j\delta_{j/b_n}$. Then
\begin{equation}\label{eq:measurelimit}
 \int(e^{\ii tx}-1-\ii tx)\nu_n(\dd x)\longrightarrow J_\alpha(t)
\end{equation}
uniformly for $t$ in compact sets. If $\phi_n(t)$ is the exponential of
the left side, its restriction to $[-\pi b_n,\pi b_n]$ has uniformly
integrable tails. More explicitly, for sufficiently small fixed $\eta>0$
there are $c,C,T>0$ such that in the small-frequency region, for
$T\le |t|\le\theta_0 b_n$,
\begin{equation}\label{eq:damping}
 |\phi_n(t)|\le C e^{-c|t|^{\alpha-\eta}},
\end{equation}
and the remaining fixed-angle region contributes $O(b_n e^{-cn})$.
The same estimates hold after retaining only jumps $j\le M_n$, uniformly
when $M_n/b_n$ lies in a compact subset of $(0,\infty)$.
\end{lemma}
\begin{proof}
Vague convergence $\nu_n\to\nu_\alpha$ follows from regular variation
and Riemann sums. For
$0<\eta<\min(2-\alpha,\alpha-1)$, Potter bounds give uniformly
\[
 \int_{(0,r]}x^2\nu_n(\dd x)\le C r^{2-\alpha-\eta}+o(1),\qquad
 \int_{[R,\infty)}x\nu_n(\dd x)\le C R^{1-\alpha+\eta}.
\]
The initial finite prefix contributes $O(n/b_n^2)=o(1)$.
These inequalities and the bounds on the compensated exponential used
in \cref{lem:abelian} prove \eqref{eq:measurelimit}.

For damping, write $\theta=t/b_n$. The negative logarithm of the modulus
is $n\beta_n\sum a_j(1-\cos(j\theta))$. For small nonzero $\theta$,
restrict $j$ to $1/|\theta|\le j\le2/|\theta|$. On this interval
$1-\cos(j\theta)$ is bounded below by a positive constant. The sum is at
least $c|\theta|^\alpha\ell(1/|\theta|)$. Dividing by the normalization at
$b_n$ and applying a lower Potter bound gives $c|t|^{\alpha-\eta}$ when
$1/|\theta|$ is sufficiently large. For all remaining angles bounded away
from zero, $a_1>0$ gives
$n\beta_n a_1(1-\cos\theta)\ge cn$ on $[-\pi,\pi]$.

For a cutoff with $M_n/b_n\ge s_->0$, the same block of indices is retained
once $|t|\ge2/s_-$. The finite interval of smaller $t$ does not affect
tail integrability. The $j=1$ term is retained eventually as well.
\end{proof}

\begin{theorem}[Uniform stable coefficient law]\label{thm:local}
Under \eqref{eq:assumptions} and \eqref{eq:bscale},
\begin{equation}\label{eq:llt}
 \sup_{k\in\mathbb Z}
 \left|b_n\Prob(K_n=k)
 -g_\alpha\!\left(\frac{k-n\beta_n d_1}{b_n}\right)\right|\to0.
\end{equation}
In particular, in every bounded coupling window,
\begin{equation}\label{eq:coeffleading}
 \boxed{\quad u_n(\beta_n)=\frac{\rho_{\beta_n}^{-n}}{n b_n}
 \bigl(g_\alpha(\lambda_n)+o(1)\bigr).\quad}
\end{equation}
The error is relative $o(1)$ uniformly when $\lambda_n$ stays in a compact
set.
\end{theorem}
\begin{proof}
Fourier inversion on the integer lattice gives
\[
 b_n\Prob(K_n=k)=\frac1{2\pi}\int_{-\pi b_n}^{\pi b_n}
 e^{-\ii t(k-n\beta_n d_1)/b_n}\phi_n(t)\dd t.
\]
By \cref{lem:fourier}, the integrands without the unit-modulus phase
converge in $L^1(\R)$ after extension by zero. This proves the supremum
statement independently of $k$. The stable density is strictly positive
on $\R$; a direct positivity argument is included in the next section.
Its positive minimum on each compact set upgrades the result to relative
error there. Equation \eqref{eq:poisson} completes the proof.
\end{proof}

The exponential factor $\rho_{\beta_n}^{-n}$ must not be replaced by
$\rhoc^{-n}$ merely because $\beta_n\to\bc$. Their ratio can be exponentially
large or small on the critical scale. Likewise, $\lambda_n$ is the actual
normalized mean defect, not a coordinate that may be dropped prematurely.

\section{The universal cutoff profile and the whole large-action cloud}
\label{sec:cutoff}

For $s>0$ define
\begin{align}
 \Psi_{\alpha,s}(t)
 &=J_\alpha(t)-\int_s^\infty e^{\ii tx}x^{-1-\alpha}\dd x,
 \label{eq:Psi}\\
 H_{\alpha,s}(y)&=\frac1{2\pi}\int_\R
                    e^{-\ii ty+\Psi_{\alpha,s}(t)}\dd t,
 \qquad
 \Phi_\alpha(s,y)=\frac{H_{\alpha,s}(y)}{g_\alpha(y)}.
 \label{eq:Phi}
\end{align}
The function $H$ is a subprobability density. The uncompensated integral
removed in \eqref{eq:Psi} is essential: replacing it by a compensated
integral would incorrectly renormalize the truncated model.

An equivalent useful expression is
\begin{equation}\label{eq:Hs}
 H_{\alpha,s}(y)=e^{-s^{-\alpha}/\alpha}
 g_{\alpha,s}\!\left(y+\frac{s^{1-\alpha}}{\alpha-1}\right),
\end{equation}
where $g_{\alpha,s}$ is the probability density with characteristic
exponent $\int_0^s(e^{\ii tx}-1-\ii tx)x^{-1-\alpha}\dd x$.
This separates the probability of no large jump from the mean shift caused
by removing those jumps.
% ed. (2026-09-29): normalization dictionary with the critical Hahn article, checked on filing (batch 50).
\begin{ednote}
In the notation of the critical Hahn article \cite{repo-critical},
$b_n^{\mathrm{CHT}}=\Gamma(-\alpha)^{1/\alpha}b_n$,
$L=s\,\Gamma(-\alpha)^{-1/\alpha}$,
$\lambda^{\mathrm{CHT}}=\Gamma(-\alpha)^{-1/\alpha}\lambda_n$, and
$\Phi_\alpha(s,y)$ is its profile $R_\alpha$ at those arguments. With
$\ell\equiv1$, \cref{thm:cutoff,thm:budget} are its \texttt{thm:cutoff}
and \texttt{thm:budget} (checked at $\alpha=3/2$, $s=0.99975$: both give
$0.357773962623$, the value recorded in \texttt{verification/cutoffs.csv}).
\end{ednote}

\begin{theorem}[Universal coupling--cutoff law]\label{thm:cutoff}
If $s_n=M_n/b_n$ remains in a compact subset of $(0,\infty)$, then
\begin{equation}\label{eq:cutofflimit}
 \boxed{\quad
 \frac{u_{n,M_n}(\beta_n)}{u_n(\beta_n)}
 =\Phi_\alpha(s_n,\lambda_n)+o(1),\quad}
\end{equation}
uniformly on compact sets of $(s_n,\lambda_n)$. No slowly varying factor
appears in the profile itself; all its leading influence is in $b_n$,
$\beta_n$, and $\rho_{\beta_n}$.
\end{theorem}
\begin{proof}
Let $B_{n,M}$ be the no-jump-above-$M$ event. Its centered Fourier transform is
\[
 \E\left[e^{\ii t(K_n-n\beta_n d_1)/b_n}\mathbf1_{B_{n,M}}\right]
 =\exp\left\{\int(e^{\ii tx}-1-\ii tx)\nu_n(\dd x)
              -\int_{(M/b_n,\infty)}e^{\ii tx}\nu_n(\dd x)\right\}.
\]
The exponent converges to \eqref{eq:Psi}. Convergence is uniform for
$s=M/b_n$ in compact subsets of $(0,\infty)$: regular variation controls
the tails, and the limiting measure has no atom at a moving cutoff.
The truncated damping in \cref{lem:fourier} gives $L^1$ convergence of the
Fourier integrands. Consequently
\[
 b_n\Prob(K_n=n,B_{n,M_n})=H_{\alpha,s_n}(\lambda_n)+o(1).
\]
Divide by \eqref{eq:llt} and use \eqref{eq:ratio}.
\end{proof}

\begin{proposition}[Positivity, a derivative identity, and endpoint values]
\label{prop:profile}
For all $s>0$ and $y\in\R$,
\begin{equation}\label{eq:strict}
 0<\Phi_\alpha(s,y)<1,
 \qquad
 \partial_s\Phi_\alpha(s,y)
 =s^{-1-\alpha}\frac{H_{\alpha,s}(y-s)}{g_\alpha(y)}>0.
\end{equation}
Moreover
\begin{equation}\label{eq:ends}
 \Phi_\alpha(s,y)\to0\quad(s\downarrow0),\qquad
 \Phi_\alpha(s,y)\to1\quad(s\to\infty),
\end{equation}
locally uniformly in $y$. All these functions are smooth in $s,y,\alpha$
for $s>0$ and $1<\alpha<2$.
\end{proposition}
\begin{proof}
The bounded-jump characteristic function has decay $e^{-c|t|^\alpha}$
for large $|t|$, by the continuum version of the damping proof. Its density
therefore extends to an entire function of $y$. The same is true with half
the L\'evy intensity. That latter density is nonnegative, not identically
zero, and real analytic, so its positive set is open and dense on the real
line. The convolution of this density with itself is consequently strictly
positive at every real point. This proves $g_{\alpha,s}(y)>0$. The full
stable law is its convolution with the independent compensated large-jump
law, so $g_\alpha(y)>0$ too. Equation \eqref{eq:Hs} gives $H>0$.

The full density decomposes according to the number of jumps above $s$.
Its zero-jump term is $H_{\alpha,s}(y)$; its one-jump term is
$\int_s^\infty H_{\alpha,s}(y-x)x^{-1-\alpha}\dd x$, which is strictly
positive. Hence $H<g$. Differentiating \eqref{eq:Psi} gives
$\partial_s\Psi_{\alpha,s}(t)=e^{\ii ts}s^{-1-\alpha}$. Differentiation
under the Fourier integral is justified locally by its exponential
majorant and gives \eqref{eq:strict}.

As $s\to\infty$, the removed integral in \eqref{eq:Psi} has absolute value
at most $s^{-\alpha}/\alpha$. Dominated Fourier convergence proves $H\to g$.
For the other endpoint, exponentially tilt the bounded-jump density by
$z=c/s$, with a sufficiently small fixed $c>0$. Its tilted characteristic
function has modulus bounded by the untilted one, since its jump intensity
is multiplied by $e^{zx}\ge1$. Splitting its Fourier integral at $1/s$
therefore bounds its density by $C/s$ for $s\le1$. With
\[
 I(c)=\int_0^1(e^{cv}-1-cv)v^{-1-\alpha}\dd v=O(c^2),
\]
\eqref{eq:Hs} yields, uniformly for $y$ in a compact set,
\[
 H_{\alpha,s}(y)\le\frac Cs
 \exp\left\{s^{-\alpha}\left(-\frac1\alpha-\frac c{\alpha-1}+I(c)\right)
                 -\frac{cy}s\right\}.
\]
The coefficient of $s^{-\alpha}$ is negative for small $c$. Because
$\alpha>1$, this dominates the $O(s^{-1})$ term and proves $H\to0$.
Local uniformity and smoothness follow from the same Fourier bounds,
including logarithmic factors generated by $\alpha$ derivatives.
\end{proof}

\begin{theorem}[Sharp action budget]\label{thm:budget}
For any integer sequence $M_n\ge1$ in a bounded coupling window,
\begin{equation}\label{eq:budgetiff}
 \frac{u_{n,M_n}(\beta_n)}{u_n(\beta_n)}\longrightarrow1
 \quad\Longleftrightarrow\quad \frac{M_n}{b_n}\longrightarrow\infty.
\end{equation}
If $M_n/b_n\to0$, the retained fraction tends to zero. If
$M_n/b_n\to s\in(0,\infty)$ and $\lambda_n\to\lambda$, it tends to
$\Phi_\alpha(s,\lambda)\in(0,1)$.
\end{theorem}
\begin{proof}
The finite positive limit follows from \cref{thm:cutoff}. If $M_n/b_n\to
\infty$, compare from below with $\lfloor s b_n\rfloor$ for a fixed large
$s$, use monotonicity from \cref{prop:exact}, and then let $s\to\infty$ in
\cref{prop:profile}. The same comparison from above with a fixed small $s$
proves the zero-budget statement. If $M_n/b_n$ does not tend to infinity,
a subsequence has bounded ratios; take a further subsequence on which
those ratios and $\lambda_n$ converge. The zero and finite positive cases
give a limit strictly below one, a contradiction.
\end{proof}

\subsection{A marked point-process formulation}
The maximum action is only one observable. Let
$\Xi_n=\sum_jN_{n,j}\delta_{j/b_n}$. For a continuous $f\ge0$ with compact
support in $(0,\infty)$, put
\begin{align}
 \Psi_{\alpha,f}(t)&=J_\alpha(t)
 +\int_0^\infty(e^{-f(x)}-1)e^{\ii tx}x^{-1-\alpha}\dd x,\label{eq:markexp}\\
 H_{\alpha,f}(y)&=\frac1{2\pi}\int_\R
                   e^{-\ii ty+\Psi_{\alpha,f}(t)}\dd t.\label{eq:markdensity}
\end{align}

\begin{theorem}[Conditioned large-action cloud]\label{thm:cloud}
If $\lambda_n\to\lambda$, then
\begin{equation}\label{eq:cloud}
 \E\left[e^{-\langle f,\Xi_n\rangle}\mid K_n=n\right]
 \longrightarrow\frac{H_{\alpha,f}(\lambda)}{g_\alpha(\lambda)}.
\end{equation}
The right side is the Laplace functional of the Poisson jump measure
$\nu_\alpha$, conditioned on its compensated stable total being $\lambda$,
with conditioning defined through these continuous densities.
\end{theorem}
\begin{proof}
Poisson marking multiplies the intensity at jump $j/b_n$ by
$e^{-f(j/b_n)}$ in the exponential formula, giving exactly the discrete
version of \eqref{eq:markexp}. The extra integral is over a compact set away
from zero. All sufficiently small jumps remain unmarked, so the Fourier
damping proof is unchanged. Fourier inversion and division by the local
limit prove \eqref{eq:cloud}. To identify the limit, split the Poisson
measure at any level below the support of $f$. Its small-jump component
has a strictly positive density; integrating the independent marked
large-jump law against that density gives precisely \eqref{eq:markdensity}.
This also defines the conditional law at every real $\lambda$, not just
almost everywhere. Standard Laplace-functional characterization then
identifies the locally finite point-process limit away from zero.
\end{proof}

This is a conditioned, not an unconditioned, Poisson cloud. At $1<\alpha<2$
the large jumps live on the same scale as the stable total, so conditioning
does not become negligible. In particular the retained fraction is generally
not the elementary no-jump factor $e^{-s^{-\alpha}/\alpha}$.

\section{Exact logarithmic tails and a convergent Lambert chart}
\label{sec:chart}

From this point until \cref{sec:computations}, strengthen the tail assumption
to
\begin{equation}\label{eq:logmodel}
 A(t)=a(-\partial_\alpha)^m\Li_{1+\alpha}(t)+P(t),
 \qquad m\in\{1,2,\ldots\},
\end{equation}
where $P$ is a real polynomial with zero constant term, and the resulting
$a_j$ are nonnegative with $a_1>0$. Thus the eventual tail is exactly
$a j^{-1-\alpha}(\log j)^m$. The polynomial changes any finite prefix.
Its coefficients need not themselves be nonnegative.

Define
\begin{equation}\label{eq:dklog}
 d_k=a(-1)^m\zeta^{(m)}(1+\alpha-k)+\sum_j [t^j]P(t)\,j^k.
\end{equation}
For $k=0,1$ these agree with \eqref{eq:constants}. For $k\ge2$ they are
analytic expansion coefficients, not convergent moments.
Parameter differentiation of the convergent polylogarithm identity
\cite[Eq.~25.12.12]{dlmf-polylog} gives, for $|x|<2\pi$ on a slit branch,
\begin{equation}\label{eq:lognormal}
 A(e^{-x})=d_0-d_1x+a\Gamma(-\alpha)x^\alpha Q_m(\log(1/x))
                  +\sum_{k\ge2}\frac{d_k(-x)^k}{k!},
\end{equation}
where
\begin{equation}\label{eq:Q}
 Q_m(L)=\sum_{r=0}^m p_r L^{m-r},\qquad
 p_r=\binom mr\frac{\Gamma^{(r)}(-\alpha)}{\Gamma(-\alpha)},\qquad p_0=1.
\end{equation}
For example, $p_1=m\psi(-\alpha)$ and
$p_2=\binom m2(\psi(-\alpha)^2+\psi_1(-\alpha))$.
The sign of $p_1$ matters; the differentiation operator is
$-\partial_\alpha$, not $\partial_\alpha$.

At $\beta=\bc$, put $C=\bc a\Gamma(-\alpha)$ and
\[
 H(x)=\bc\sum_{k\ge2}\frac{d_k(-1)^k x^{k-2}}{k!}.
\]
The exact critical normal form is
\begin{equation}\label{eq:exactnormal}
 \boxed{\quad\varepsilon=Cx^\alpha Q_m(\log(1/x))+x^2H(x).\quad}
\end{equation}
This is an identity of analytic functions on the slit neighborhood, not
merely an asymptotic expansion.

\subsection{A Lambert core that keeps the leading logarithm exact}
Let $T=e^{-L}$ be the small positive solution of
\begin{equation}\label{eq:core}
 T^\alpha L^m=\varepsilon/C,\qquad L>m/\alpha.
\end{equation}
Its exact expression is
\begin{equation}\label{eq:Lambertcore}
 L=-\frac m\alpha W_{-1}\!\left(
      -\frac\alpha m(\varepsilon/C)^{1/m}\right),\qquad T=e^{-L}.
\end{equation}
The $W_{-1}$ branch is forced by $L\to\infty$ as $\varepsilon\to0$.
The other real branch would not describe the critical boundary.

\begin{theorem}[Absolutely convergent three-coordinate inverse chart]
\label{thm:chart}
There is a holomorphic function $V(u,T,R)$ near $(0,0,0)$ with
$V(0,0,0)=1$ such that the positive critical inverse, and its local complex
continuation on compatible branches, is exactly
\begin{equation}\label{eq:chart}
 x=T V\left(\frac1L,T,\frac{T^{2-\alpha}}{L^m}\right).
\end{equation}
It is characterized by
\begin{equation}\label{eq:Veq}
 V^\alpha\sum_{r=0}^m p_r u^r(1-u\log V)^{m-r}
       +\frac RC V^2H(TV)=1.
\end{equation}
Its Taylor series converges absolutely on a sufficiently small polydisc.
After substitution, its monomials are contained in the finite-generator
scale
\begin{equation}\label{eq:chartgrid}
 \left\{T^{1+j+(2-\alpha)k}L^{-i-mk}:i,j,k\in\N_0\right\}.
\end{equation}
There is no finite Puiseux chart in $\varepsilon$ for $m\ge1$.
\end{theorem}
\begin{proof}
Substitute $x=TV$ into \eqref{eq:exactnormal} and divide by
$C T^\alpha L^m$. Since $\log(1/x)=L-\log V$, the result is
\eqref{eq:Veq} after setting $u=1/L$ and $R=T^{2-\alpha}L^{-m}$.
At $(u,T,R,V)=(0,0,0,1)$ the left side minus one vanishes and its derivative
in $V$ is $\alpha\ne0$. The holomorphic implicit-function theorem yields
$V$. Its Taylor series has the usual absolute convergence on smaller
polydiscs, and substitution yields \eqref{eq:chartgrid}.

These are a finite set of small generators over the exact Lambert core.
At a fixed power of $T$, there may be an entire convergent inverse-logarithmic
block; the claim is not local finiteness with respect to the $T$ exponent
alone. In the lexicographically ordered power--log scale the support is
well based. If $\alpha$ is irrational, the positive $T$-power generators
still have only finitely many pairs $(j,k)$ below any fixed bound.

Finally, $x\sim T$ and
$T=\varepsilon^{1/\alpha}$ times a nonconstant logarithmic factor tending
to zero. Thus $x$ is not asymptotic to a constant times a rational power
of $\varepsilon$, excluding a finite Puiseux chart even when $\alpha$
is rational.
\end{proof}

\begin{proposition}[First chart coefficients]\label{prop:chartcoeff}
With $p_2=0$ for $m=1$,
\begin{align}
 V(u,0,0)
 &=1-\frac{p_1}{\alpha}u
 +\left(\frac{(\alpha+1)p_1^2}{2\alpha^2}
        -\frac{mp_1}{\alpha^2}-\frac{p_2}{\alpha}\right)u^2+O(u^3),
 \label{eq:Vtwo}\\
 [R]V(0,0,R)&=-\frac{H(0)}{\alpha C}.\label{eq:VR}
\end{align}
\end{proposition}
\begin{proof}
Write $V=1+v_1u+v_2u^2+\cdots$. Up to order $u^2$,
\[
 \sum_{r=0}^m p_r u^r(1-u\log V)^{m-r}
 =1+p_1u+(p_2-mv_1)u^2+O(u^3).
\]
Multiplying by the binomial expansion of $V^\alpha$ and setting the first
two coefficients equal to zero gives \eqref{eq:Vtwo}. Differentiation of
\eqref{eq:Veq} in $R$ at the origin gives \eqref{eq:VR}.
\end{proof}

\subsection{A genuine geometric error bound}
Suppose $V$ is bounded by $K_0$ on a polydisc with radii $u_0,T_0,R_0$,
and let $V_{\le N}$ be its Taylor polynomial of total degree at most $N$.
At the physical coordinates put
\[
 \theta=\max\left\{\frac1{|L| u_0},\frac{|T|}{T_0},
              \frac{|T|^{2-\alpha}|L|^{-m}}{R_0}\right\}<1.
\]
Cauchy's estimate gives
\begin{equation}\label{eq:geometric}
 \left|x-T V_{\le N}(L^{-1},T,T^{2-\alpha}L^{-m})\right|
 \le |T|K_0\binom{N+3}{2}\frac{\theta^{N+1}}{(1-\theta)^3}.
\end{equation}
Indeed, degree $k$ contains $\binom{k+2}{2}$ three-variable monomials;
summing their geometric majorant proves the bound.

The radii can be chosen by a finite analytic inequality. On
$|v-1|=1/4$, let $m_\alpha=\min|v^\alpha-1|>0$, with powers and logarithms
on their branches near one. Choose $T_0$ so $|Tv|$ remains inside a disc
of analyticity of $H$, and choose $u_0,R_0$ so the difference between the
left side of \eqref{eq:Veq} and $v^\alpha$ has modulus less than
$m_\alpha/2$ on that circle. Rouch\'e's theorem gives a unique simple zero
inside it, uniformly in the parameters. Slightly shrinking the radii then
makes $K_0=5/4$ a valid Cauchy bound. Numerical evaluation with directed
rounding would turn these inequalities into explicit interval certificates;
the supplied diagnostic program does not claim to do that step.

\section{All-order logarithmic coefficient and cutoff expansions}
\label{sec:allorders}

For the exact logarithmic model, choose the scale by the exact equation
\begin{equation}\label{eq:exactscale}
 b_n^\alpha=n\beta_n a(\log b_n)^m,\qquad b_n>e^{m/\alpha},
 \qquad L_n=\log b_n.
\end{equation}
The large solution is unique for sufficiently large $n$ and is
\begin{equation}\label{eq:bnLambert}
 L_n=-\frac m\alpha W_{-1}\!\left(
    -\frac\alpha m(n\beta_n a)^{-1/m}\right).
\end{equation}
In particular
\begin{equation}\label{eq:bnleading}
 b_n\sim(n\beta_n a)^{1/\alpha}
       \left(\frac{\log(n\beta_n a)}\alpha\right)^{m/\alpha}.
\end{equation}
Replacing \eqref{eq:exactscale} by \eqref{eq:bnleading} is valid at leading
order but generally invalid for the correction formulas below.

For either $Q_\alpha(t)=J_\alpha(t)$ or
$Q_\alpha(t)=\Psi_{\alpha,s}(t)$, define
\begin{equation}\label{eq:Bbell}
 B_k[Q](t)=
 \sum_{\substack{r_1,\ldots,r_m\ge0\\\sum_{h=1}^m h r_h=k}}
 \prod_{h=1}^m\frac{\left(\binom mh(-\partial_\alpha)^h Q_\alpha(t)
                      \right)^{r_h}}{r_h!},
 \qquad B_0[Q]=1.
\end{equation}
Derivatives in this section always hold $t$, $s$, and the normalized
coordinate $y$ fixed. They do not differentiate $b_n$, $\beta_n$, $d_1$,
or $\lambda_n$.
Write
\begin{align}
 D_k(y)&=\frac1{2\pi}\int_\R e^{-\ii ty+J_\alpha(t)}B_k[J](t)\dd t,
 \label{eq:Dk}\\
 N_k(s,y)&=\frac1{2\pi}\int_\R
                 e^{-\ii ty+\Psi_{\alpha,s}(t)}B_k[\Psi_s](t)\dd t.
 \label{eq:Nk}
\end{align}
The letter $D_k$ here denotes a correction function, not the analytic
coefficient $d_k$ in \eqref{eq:dklog}. In particular $D_0=g_\alpha$ and
$N_0=H_{\alpha,s}$.

\begin{theorem}[All fixed orders]\label{thm:allorders}
For every fixed nonnegative integer $K$, the exact logarithmic model has
\begin{align}
 n b_n\rho_{\beta_n}^{n}u_n(\beta_n)
 &=\sum_{k=0}^K\frac{D_k(\lambda_n)}{L_n^k}+O(L_n^{-K-1}),
 \label{eq:allfull}\\
 n b_n\rho_{\beta_n}^{n}u_{n,M_n}(\beta_n)
 &=\sum_{k=0}^K\frac{N_k(s_n,\lambda_n)}{L_n^k}+O(L_n^{-K-1}),
 \label{eq:allcut}
\end{align}
uniformly for $s_n=M_n/b_n$ and $\lambda_n$ in compact subsets of
$(0,\infty)$ and $\R$, respectively. Finite-prefix changes enter the
normalization and algebraically smaller corrections but not the displayed
coefficient functions.
\end{theorem}
\begin{proof}
We give the error analysis, since a fixed-frequency formal expansion alone
would not prove the theorem. Parameter differentiation of
\eqref{eq:lognormal}, with $x=-\ii t/b_n$, gives the centered exponent
\begin{equation}\label{eq:exponentexp}
 \log\phi_n(t)=J_\alpha(t)
 +\sum_{h=1}^m\binom mh L_n^{-h}(-\partial_\alpha)^hJ_\alpha(t)
 +E_n(t),
\end{equation}
where, for $|t|\le\theta_0b_n$,
\begin{equation}\label{eq:analyticerror}
 |E_n(t)|\le C\frac n{b_n^2}|t|^2.
\end{equation}
To see the powers of $L_n$, note that the singular part is
$(-\partial_\alpha)^m\{b_n^{-\alpha}J_\alpha(t)\}$ with $b_n$ held
fixed during this algebraic differentiation; Leibniz's rule and
\eqref{eq:exactscale} give \eqref{eq:exponentexp}.
The analytic remainder follows from the convergent analytic series in
\eqref{eq:lognormal}. Since $\alpha<2$,
\[
 \frac n{b_n^2}=\frac{b_n^{\alpha-2}}{\beta_n a L_n^m}
              =o(L_n^{-r})\quad\text{for every fixed }r.
\]

For the cutoff exponent, subtract
$n\beta_n\sum_{j>M_n}a_je^{\ii tj/b_n}$. For large $n$ all indices here
are beyond the finite prefix. The normalized weights on the mesh $j/b_n$
are exactly
\[
 \frac1{b_n}(j/b_n)^{-1-\alpha}
       \left(1+\frac{\log(j/b_n)}{L_n}\right)^m.
\]
On $s_n\ge s_->0$, comparison with its integral over $(s_n,\infty)$
gives an error $O((1+|t|)/b_n)$. This follows by bounding the integral of
the absolute derivative of
$e^{\ii tx}x^{-1-\alpha}(1+\log x/L_n)^m$; its logarithmic moments at
infinity are finite, uniformly for $L_n\ge1$. Expanding the finite
polynomial in $L_n^{-1}$ shows that \eqref{eq:exponentexp} becomes
\begin{equation}\label{eq:cutexponentexp}
 \Psi_{\alpha,s_n}(t)+
 \sum_{h=1}^m\binom mhL_n^{-h}(-\partial_\alpha)^h\Psi_{\alpha,s_n}(t)
 +O\left(\frac n{b_n^2}t^2+\frac{1+|t|}{b_n}\right).
\end{equation}
The identity $(-\partial_\alpha)^h x^{-1-\alpha}
=(\log x)^h x^{-1-\alpha}$ proves that the same differential operator
acts on both the retained and removed parts, including the killing and
centering terms.

Choose a small $\kappa>0$, for example $\kappa<1/(4\alpha)$, and split the
Fourier integral at $|t|=L_n^\kappa$. Beyond this interval,
\cref{lem:fourier} bounds the actual Fourier tails by
$\exp(-cL_n^{\kappa(\alpha-\eta)})$, smaller than every inverse power of
$L_n$. The tails of every proposed coefficient integral have the same
property. On $1\le|t|\le L_n^\kappa$, the perturbation in
\eqref{eq:exponentexp} is $o(1)$ relative to the magnitude of the negative
real part of $J_\alpha(t)$. For the truncated exponent, the difference
from $J_\alpha(t)$ and its fixed $\alpha$ derivatives is bounded uniformly
in $s$ on compact sets; the same conclusion holds for sufficiently large
$|t|$. On $|t|\le1$, the quantities
$|t|^\alpha(1+|\log|t||^m)$ are bounded, so the perturbations are $O(1/L_n)$
uniformly. The analytic errors in \eqref{eq:analyticerror} and
\eqref{eq:cutexponentexp} are smaller than every inverse-logarithmic order
on the retained frequency interval.

Taylor expansion of the exponential through total order $K$ now has an
integrable remainder bounded by $C_K L_n^{-K-1}$ times a polynomial in
$|t|$ and $|\log|t||$, multiplied by $e^{-c|t|^\alpha}$ outside a fixed
interval. Its coefficients are exactly \eqref{eq:Bbell}. Integration
proves \eqref{eq:allfull} and \eqref{eq:allcut}. The phase
$e^{-\ii t\lambda_n}$ has unit modulus, so compact uniformity in
$\lambda_n$ presents no further error. This also explains why finite
prefixes contribute only below all fixed inverse-logarithmic orders once
their effect on $d_0,d_1,\bc$ has been retained.
\end{proof}

\begin{remark}[Uniformity in the exponent]
The same argument is uniform for $\alpha$ in any fixed compact subinterval
of $(1,2)$, for a correspondingly controlled family of prefixes. It is not
uniform as $\alpha$ approaches either endpoint. The factors controlling
small jumps, large compensated jumps, and $n/b_n^2$ degenerate there.
\end{remark}

\subsection{The effective-index law}
Since $B_1[Q]=-m\partial_\alpha Q$, differentiation under the Fourier
integral gives
\begin{equation}\label{eq:firstDN}
 D_1(y)=-m\partial_\alpha g_\alpha(y),\qquad
 N_1(s,y)=-m\partial_\alpha H_{\alpha,s}(y).
\end{equation}
Dividing the two expansions yields the next result.

\begin{corollary}[A common first correction]\label{cor:effective}
For the exact logarithmic model,
\begin{align}
 n b_n\rho_{\beta_n}^{n}u_n(\beta_n)
 &=g_\alpha(\lambda_n)-\frac m{L_n}\partial_\alpha g_\alpha(\lambda_n)
   +O(L_n^{-2}),\label{eq:firstfull}\\
 \frac{u_{n,M_n}(\beta_n)}{u_n(\beta_n)}
 &=\Phi_\alpha(s_n,\lambda_n)
   -\frac m{L_n}\partial_\alpha\Phi_\alpha(s_n,\lambda_n)
   +O(L_n^{-2}).\label{eq:firstPhi}
\end{align}
Equivalently the two leading profiles may be evaluated at
\begin{equation}\label{eq:effectivealpha}
 \alpha_{\mathrm{eff},n}=\alpha-\frac m{\log b_n},
\end{equation}
with an $O((\log b_n)^{-2})$ error, provided $b_n$, $s_n$, and $\lambda_n$
are held fixed in this evaluation.
\end{corollary}
\begin{proof}
The quotient expansion gives
\[
 \frac{-m\partial_\alpha H}{g}
 -\frac{H(-m\partial_\alpha g)}{g^2}
 =-m\partial_\alpha(H/g).
\]
Taylor's theorem in $\alpha$ proves the last interpretation.
\end{proof}

The effective index has a direct interpretation at the intensity level:
\[
 x^{-1-\alpha}\left(1+\frac{\log x}{L_n}\right)^m
 =x^{-1-\alpha}\left(1+\frac m{L_n}\log x+O(L_n^{-2}\log^2x)\right).
\]
Its first perturbation is the same as that of
$x^{-1-(\alpha-m/L_n)}$. This explanation is not a substitute for the
Fourier proof: the latter handles the nonuniform logarithms near zero
and infinity, the lattice, and the coefficient conditioning.

\subsection{Higher orders are not just an index shift}
At the next order,
\begin{equation}\label{eq:secondB}
 B_2[Q]=\binom m2\partial_\alpha^2Q
             +\frac{m^2}{2}(\partial_\alpha Q)^2.
\end{equation}
For example,
\[
 D_2=\frac{m^2}{2}\partial_\alpha^2g_\alpha
 -\frac m2\F\bigl[e^{J_\alpha}\partial_\alpha^2J_\alpha\bigr].
\]
The second term would be missed by blindly replacing $\alpha$ everywhere
by $\alpha-m/L_n$. For arbitrary order, the quotient coefficients $R_k$
of the retained fraction are computed by the finite recursion
\begin{equation}\label{eq:ratiorec}
 R_0=\Phi_\alpha,\qquad
 R_k(s,y)=\frac{N_k(s,y)-\sum_{j=1}^kD_j(y)R_{k-j}(s,y)}{g_\alpha(y)}.
\end{equation}
Thus $\sum_{k\le K}R_k L_n^{-k}$ has the same fixed-order uniform error as
in \cref{thm:allorders}. No convergence in $K$ is claimed.

\section{Two-term minimal action budgets}
\label{sec:quantiles}

Fix a desired retained fraction $p\in(0,1)$. By \cref{prop:profile}, the
equation
\begin{equation}\label{eq:quantile}
 \Phi_\alpha(s_p(\alpha,y),y)=p
\end{equation}
has exactly one solution $s_p(\alpha,y)>0$. The derivative formula
\eqref{eq:strict} and the implicit-function theorem show that $s_p$ is
smooth and
\begin{equation}\label{eq:quantilederivative}
 \partial_\alpha s_p(\alpha,y)
 =-\frac{\partial_\alpha\Phi_\alpha(s_p(\alpha,y),y)}
         {\partial_s\Phi_\alpha(s_p(\alpha,y),y)}.
\end{equation}

Let $M_n^*(p)$ be the smallest integer $M$ for which
$u_{n,M}(\beta_n)/u_n(\beta_n)\ge p$. It exists because the ratio equals
one at $M=n$.

\begin{theorem}[An explicit two-term asymptotic budget]\label{thm:minbudget}
In the exact logarithmic model and a bounded coupling window,
\begin{align}
 \frac{M_n^*(p)}{b_n}
 &=s_p(\alpha,\lambda_n)
 +\frac m{L_n}
 \frac{\partial_\alpha\Phi_\alpha(s_p(\alpha,\lambda_n),\lambda_n)}
      {\partial_s\Phi_\alpha(s_p(\alpha,\lambda_n),\lambda_n)}
 +O(L_n^{-2}+b_n^{-1})\label{eq:minbudget}\\
 &=s_p\!\left(\alpha-\frac m{L_n},\lambda_n\right)
 +O(L_n^{-2}+b_n^{-1}).\label{eq:minbudgeteff}
\end{align}
The denominator in \eqref{eq:minbudget} is explicitly positive by
\eqref{eq:strict}. The integer ambiguity is at most the indicated
$O(1)$ in $M_n^*(p)$.
\end{theorem}
\begin{proof}
For $y$ in a compact set the quantiles stay in a compact subinterval of
$(0,\infty)$, by local uniformity of \eqref{eq:ends}. On a slightly larger
compact rectangle, $\partial_s\Phi_\alpha$ has a positive lower bound.
Insert
$s=s_p+(m/L_n)(\partial_\alpha\Phi/\partial_s\Phi)$ into
\eqref{eq:firstPhi}; its order-$L_n^{-1}$ terms cancel. The remaining
error is $O(L_n^{-2})$ uniformly. Moving $s$ by a sufficiently large
constant times $L_n^{-2}$ brackets the target value by the derivative
lower bound. Monotonicity of the actual retained fraction then brackets
its smallest integer crossing, and rounding changes $s$ by at most
$1/b_n$. Equation \eqref{eq:quantilederivative} gives the second form.
\end{proof}

This is an asymptotic selection rule, not an unconditional finite-$n$
certificate. A theorem with an unspecified $O(L_n^{-2})$ constant cannot
by itself promise a particular retained fraction at a stated finite $n$.
A proof-producing selector would additionally enclose the Fourier integrals,
provide effective constants for the error analysis, and round outward.

\section{Reproducible calculations and what they actually verify}
\label{sec:computations}

\subsection{Exact checks}
The supplied \texttt{verification/verify.py} performs 552 exact assertions.
It compares direct feedback recursion with Lagrange extraction for three
rational couplings through order 16, checks finite-cutoff monotonicity and
agreement through the cutoff, verifies the first two chart coefficients
symbolically, and checks a Bell-polynomial coefficient recurrence through
degree six. These are finite exact algebra checks. They do not certify
an infinite Fourier integral or an asymptotic remainder.

\subsection{A logarithmic-tail example}
The numerical example is
\begin{equation}\label{eq:numericalmodel}
 \alpha=\tfrac32,\qquad a=1,\qquad m=1,\qquad
 A(t)=t-\left.\partial_s\Li_s(t)\right|_{s=5/2}.
\end{equation}
Thus $a_1=1$ and $a_j=(\log j)/j^{5/2}$ for $j\ge2$. The constants are
\begin{align*}
 d_0&=1-\zeta'(5/2)=1.38734195032620997271\ldots,\\
 d_1&=1-\zeta'(3/2)=4.93223973743110151071\ldots,\\
 \bc&=0.20274764675588095470\ldots.
\end{align*}
At the critical coupling, $\lambda_n=0$. The first correction constants are
\begin{align*}
 D_0(0)&= 0.14026098237322958073\ldots,\\
 D_1(0)&=-0.22110216101538913083\ldots,\\
 D_2(0)&=-0.13928691395372795423\ldots,\\
 D_3(0)&= 0.70330022056722715680\ldots,\\
 D_4(0)&= 0.48955185500793247479\ldots.
\end{align*}
For $m=1$ these are computed from
\[
 D_k(0)=\frac1{2\pi k!}\int_\R e^{J_\alpha(t)}
 \left[J_\alpha(t)\bigl(\psi(-\alpha)-\log(-\ii t)\bigr)\right]^k\dd t.
\]
The expression at $t=0$ is interpreted by continuity.

\begin{table}[htbp]\centering\small
\begin{tabular}{rrrrr}\toprule
$n$ & $b_n$ & Observed $b_n\Prob(K_n=n)$ & Through $L_n^{-1}$ & Through $L_n^{-4}$\\\midrule
256 & 31.834 & 0.082557036 & 0.076368579 & 0.085122191 \\
1,024 & 96.559 & 0.093477009 & 0.091881350 & 0.093702738 \\
4,096 & 279.735 & 0.101079624 & 0.101015630 & 0.101046244 \\
16,384 & 788.914 & 0.106639717 & 0.107115498 & 0.106601915 \\
65,536 & 2185.549 & 0.110858081 & 0.111507663 & 0.110838856 \\
\bottomrule\end{tabular}
\caption{Coefficient diagnostics for the logarithmic tail at $\alpha=3/2$. The common leading approximation is $0.140260982\ldots$. These are floating-point diagnostics, not certified enclosures.}\label{tab:coeff}
\end{table}


The convergence is logarithmically slow. Additional terms need not improve
an approximation monotonically at modest $n$: the signs and sizes of
$D_k$ matter, and the smaller analytic and lattice corrections have not
been included in this table. The theorem is a fixed-order asymptotic law,
not a finite-$n$ alternation or enveloping theorem.

\begin{table}[htbp]\centering\small
\begin{tabular}{rrrrr}\toprule
$n$ & Actual $M/b_n$ & Observed retained fraction & $\Phi_{3/2}$ & First correction\\\midrule
1,024 & 0.994215 & 0.096370999 & 0.353610742 & 0.074240642 \\
4,096 & 0.997372 & 0.133121494 & 0.355987344 & 0.129040344 \\
16,384 & 0.998842 & 0.162779389 & 0.357092276 & 0.165295316 \\
65,536 & 0.999749 & 0.186518804 & 0.357773963 & 0.191326680 \\
\bottomrule\end{tabular}
\caption{Cutoff diagnostics with $M=\lfloor b_n\rfloor$ and $\lambda_n=0$. The first correction is $\Phi_{3/2}-(\log b_n)^{-1}\partial_\alpha\Phi_\alpha|_{\alpha=3/2}$.}\label{tab:cut}
\end{table}


The leading profile can be visibly inaccurate even at a large index.
The first derivative correction substantially improves the displayed
$s\approx1$ values. At smaller $s$ and modest $n$, a truncated correction
can even be negative, although the true probability never is. Such a
negative approximation is evidence that the selected truncation is not
useful there, not a violation of the probability identity or a certified
error estimate.

\subsection{Inverse-chart diagnostics}
For \eqref{eq:numericalmodel}, take $L$ as the independent core coordinate,
$T=e^{-L}$, and $\varepsilon=C T^{3/2}L$. Solve the exact normal form using
65-digit arithmetic and compare $x/T$ with \eqref{eq:Vtwo}. Here
$p_1=\psi(-3/2)$ and $p_2=0$.
\begin{table}[htbp]\centering\small
\begin{tabular}{rrrr}\toprule
$L$ & Exact-equation numerical $x/T$ & Through $L^{-1}$ & Through $L^{-2}$\\\midrule
8 & 0.939577755920 & 0.941403613280 & 0.940812501724 \\
12 & 0.960595968884 & 0.960935742186 & 0.960673025939 \\
20 & 0.976477564834 & 0.976561445312 & 0.976466867463 \\
35 & 0.986577920376 & 0.986606540178 & 0.986575657615 \\
\bottomrule\end{tabular}
\caption{Inverse-chart diagnostics. The displayed second-order approximation does not include the much smaller $R$-coordinate contribution.}\label{tab:inverse}
\end{table}


The exact equation is evaluated through a parameter derivative of the
polylogarithm, not through a truncated inverse series. This separates the
numerical check from the formula it is checking.

\subsection{Numerical method and independent cross-checks}
The coefficient computations use the exact finite-degree dependence of
\eqref{eq:lagrange}. A Fourier transform evaluates
$\exp(n\bc(A_n(re^{\ii\theta})-d_0))$ with $r=e^{-2/n}$ and transform
length at least $32n$. The coefficient at degree $n$, divided by $r^n$,
gives $\Prob(K_n=n)$ up to numerical quadrature and roundoff errors.
For a cutoff the constant remains the full $d_0$, in accordance with the
joint no-large-jump event. A separate positive recurrence,
\[
 P_0=e^{-n\bc d_0},\qquad
 P_k=\frac{n\bc}{k}\sum_{j=1}^k j a_j P_{k-j},
\]
checks the full coefficient at $n=256$ using 65-digit arithmetic.
The measured relative discrepancy is recorded in \texttt{results.json}.

Cutoff densities are evaluated by Fourier quadrature, with a Gauss--Jacobi
rule for the compensated small-jump integral. Repeating with 96 and 160
nodes checks numerical stability. The derivative in \eqref{eq:firstPhi}
is approximated by a centered difference in $\alpha$, with normalized
$s$ and $y$ held fixed. These calculations use ordinary floating-point
arithmetic and are not interval enclosures. All numerical tables are
generated from the retained machine-readable outputs; no empirical value
is used as a lemma in the proofs.

\section{Consequences for a transseries implementation}
\label{sec:implementation}

\subsection{Three distinct kinds of representation}
An implementation should distinguish the following objects instead of
forcing them into a single unchecked datatype. First, an asymptotic-inverse
object may specify $R^\asyminv$ and its regular-variation laws without an
explicit monomial series. This is the appropriate representation for
arbitrary $\ell$ in \cref{thm:rv-inverse}. Second, a convergent local chart
has a finite tuple of small coordinates, an analytic defining equation,
Cauchy radii, and a majorant such as \eqref{eq:geometric}. Third, a
coefficient-asymptotic object has correction operators and fixed-order
remainder statements but need not have a convergent infinite expansion.

The oscillatory example belongs to the first class. Exact logarithmic tails
belong to both the second and third. Keeping these distinctions explicit
prevents a formal closure claim from being mistaken for an analytic
realization theorem.

\subsection{A practical coefficient pipeline}
For the logarithmic model, a symbolic front end can compute the polynomial
$Q_m$ from gamma derivatives, generate the Taylor jet of \eqref{eq:Veq},
and obtain $B_k$ by the finite recurrence
\[
 B_0=1,\qquad k B_k=\sum_{r=1}^{\min(k,m)}r\binom mr
                      (-\partial_\alpha)^rQ\,B_{k-r}.
\]
A numerical back end then evaluates $D_k$ and $N_k$ with Fourier integrals.
The ratio recurrence \eqref{eq:ratiorec} and the positive derivative
\eqref{eq:strict} produce cutoff approximations and well-conditioned
quantile inversion on compact rectangles.

The physical scale must be obtained from \eqref{eq:exactscale} or the exact
Lambert formula, not from a rough $n^{1/\alpha}(\log n)^{m/\alpha}$
approximation when subleading accuracy is wanted. Similarly, all finite
prefix effects on $d_0$ and $d_1$ must be retained before applying the
universal normalized profile.

\subsection{A realistic formalization boundary}
No Lean file is included or claimed to compile. A useful formalization
sequence would begin with finite formal-power-series algebra: Lagrange
coefficient identities, the cutoff dependence, and the Bell recurrence.
The next layer is the analytic implicit-function theorem with an explicit
polydisc majorant. Those statements can be expressed without formalizing
stable probability laws.

The probabilistic layer needs independent Poisson sums, the exact generating
function, and Fourier inversion with an $L^1$ majorant. Only after those
interfaces are available should one assemble the regular-variation limits
and the triangular-array local theorem. Finally, a directed-rounding
quadrature checker could certify profile values and error bounds.
Existing repository results on formal support or finite recurrences are
useful ingredients, but none is a substitute for these analytic hypotheses.

\section{Further research questions and concrete targets}
\label{sec:research}

\paragraph{1. Second-order slow variation beyond exact logarithmic powers.}
Suppose
$\ell(bx)/\ell(b)=1+A(b)\log x+o(A(b))$ with a weighted uniform remainder.
Determine the weakest domination assumptions under which
$\Phi_\alpha-A(b)\partial_\alpha\Phi_\alpha$ remains the correct first
cutoff approximation. Pointwise second-order variation alone is not an
integrable error estimate; the small- and large-jump tails must be controlled.
The logarithmic model proves the target identity in one fully specified
class, not in this unrestricted second-order class.

\paragraph{2. Noninteger logarithmic powers and convergent charts.}
For $a_j=a j^{-1-\alpha}(\log j)^r$ with real noninteger $r$, the leading
regular-variation and cutoff theorems already apply. Determine whether the
inverse admits an absolutely convergent finite-coordinate chart, or only
an asymptotic chart, under a suitable analytic continuation of the tail.
Integer parameter differentiation is no longer an exact finite operation.

\paragraph{3. Effective finite-$n$ error constants.}
Make the constants in \cref{thm:allorders} explicit on a prescribed compact
rectangle. Combine interval evaluation of $D_k,N_k$ with the positive
profile derivative to produce a genuinely certified $M_n^*(p)$ enclosure.
The current two-term formula is asymptotic and the supplied numerical
outputs are not outward-rounded certificates.

\paragraph{4. Expanding coupling windows.}
Find the largest growing range of $|\lambda_n|$ on which the logarithmic
expansion remains relatively accurate, and match it to one-large-jump and
exponential-tilt regimes. The uniform absolute local theorem does not
supply relative accuracy where $g_\alpha(\lambda_n)$ becomes very small.

\paragraph{5. Joint endpoint and slow-variation limits.}
Let $\alpha\uparrow2$ or $\alpha\downarrow1$ while the logarithmic tail and
$n$ vary. Derive a common normalization and an error theorem covering the
interaction of the endpoint singularity with $m/\log b_n$. Substitution
of an endpoint in \eqref{eq:effectivealpha} is not justified by this paper.
% ed. (2026-09-29): the same limit is asked by four packages of this series (batches 48 and 49); none answers it.
\begin{ednote}
Slowly varying factors at or near $\alpha=2$ are also asked about by the
logarithmic-endpoint package's Question~6, the marginal package's
Question~2, the stable--Gaussian package's Question~3 and the confluent
package's Question~4 \cite{ed:lce,ed:mct,ed:sge,ed:cct}; none is answered.
\end{ednote}

\paragraph{6. Large order of the inverse-logarithmic expansion.}
Determine the growth of $D_k,N_k$ as $k\to\infty$, the optimal truncation
index, and whether a Borel or accelerated summation reconstructs the
coefficient law. The convergent inverse chart does not imply convergence
of this outer expansion. Its complex saddle structure may encode a second
level of transseries information.

\paragraph{7. Oscillatory slow blocks as a symbolic language.}
Develop a representation that includes
$2+\sin(\log\log y)$ without falsely asserting membership in a
nonoscillatory power--log scale. A first target is an algebra of periodic
functions of $\log\log y$ with explicit asymptotic inversion rules and
error control. Decide which coefficient and cutoff computations remain
algorithmically tractable in that enlarged language.

\paragraph{8. Perturbed slopes.}
Replace the exact slopes $j$ in \eqref{eq:feedback} by $j+r_j$.
Even $r_j=o(j)$ can exponentially reweight the boundary amplitudes through
$e^{r_jU}$, so a naive claim that the present theorem is stable under every
sublinear perturbation would be false. Identify a correctly normalized
perturbation class preserving the slow tail, then determine the new chart
generators and correction operators.

\paragraph{9. Multiple marked clouds and functional limits.}
The Laplace functional in \cref{thm:cloud} determines the limiting point
measure away from zero. Develop quantitative joint laws for ordered actions,
small-jump compensation, and the induced tree or allocation structures.
Prove metric error bounds rather than only convergence of Laplace
functionals, and track their sensitivity to the coupling coordinate.

\paragraph{10. Minimal symbolic and analytic certificates.}
Implement separate checkers for a coefficient identity, a support claim,
a convergent analytic chart, and an asymptotic error bound. A robust
ProveIt integration should make the distinction visible in its interface:
a verified finite coefficient must not silently stand in for a verified
infinite asymptotic statement.

\section{Conclusion}
The slowly varying extension has a positive and a negative side. The inverse
scale and action budget remain rigid after the correct regular-variation
normalization: the same stable density and the same conditional cutoff
profile apply. But slow variation alone does not force a power--logarithmic
inverse representation. The explicit oscillatory example separates those
two conclusions.

Exact logarithmic tails occupy a useful intermediate class. They admit a
convergent Lambert-based chart and also a computable all-order expansion
for large coefficients and action retention. At first order, the common
correction is an effective shift of the stable index. The strict derivative
formula for the cutoff profile makes that correction actionable as a
minimal-budget law. This resolves the fixed-interior-exponent part of the
repository's slowly-varying-tail question, while leaving endpoint
uniformity, unrestricted analytic realization, and fully certified
finite-$n$ computation as explicit further problems.

\clearpage
\appendix
\section{Normalization and sign audit}
\label{app:audit}
The following checks summarize the convention-sensitive points.
\begin{longtable}{p{0.27\textwidth}p{0.66\textwidth}}
\toprule
Object & Convention and consequence\\\midrule
\endhead
L\'evy measure & $x^{-1-\alpha}\dd x$, with no gamma factor. Thus
$J_\alpha(t)=\Gamma(-\alpha)(-\ii t)^\alpha$.\\[2mm]
Stable scale & $n\beta a\ell(b)/b^\alpha\to1$. A scale that absorbs
$\Gamma(-\alpha)$ requires a correspondingly rescaled density and cutoff.\\[2mm]
Coupling coordinate & $\lambda_n=(n-n\beta_nd_1)/b_n$; positive values
are targets above the mean.\\[2mm]
No-large-jump exponent & $\Psi_s=J-\int_s^\infty e^{\ii tx}\nu(\dd x)$.
The removed term is not compensated.\\[2mm]
Mean shift after removal & $H_s(y)=e^{-s^{-\alpha}/\alpha}
 g_s(y+s^{1-\alpha}/(\alpha-1))$. The shift has a plus sign inside $g_s$.\\[2mm]
Logarithmic tail & $(-\partial_\alpha)^m\Li_{1+\alpha}$, giving
$(\log j)^m j^{-1-\alpha}$.\\[2mm]
First profile correction & $-(m/\log b)\partial_\alpha\Phi_\alpha$ at
fixed normalized coordinates.\\[2mm]
Quantile correction & $+(m/\log b)(\partial_\alpha\Phi/\partial_s\Phi)$,
or equivalently $s_p(\alpha-m/\log b,y)$.\\[2mm]
Critical inverse & $\varepsilon=x+\bc(A(e^{-x})-d_0)$;
$U=\bc d_0+\varepsilon-x$.\\[2mm]
Lambert branch & $W_{-1}$, since both the inverse-core $L$ and the
coefficient-scale $\log b$ tend to $+\infty$.\\
\bottomrule
\end{longtable}

\section{Reproduction instructions}
\label{app:reproduce}
The package is self-contained apart from standard LaTeX packages and the
Python dependencies listed in \texttt{requirements.txt}. From the package
root, run
% ed. (2026-09-29): the delivered recipe rewrote the recorded outputs; the programs now write into rerun/ (see the package README).
\begin{verbatim}
python verification/verify.py
python verification/make_tables.py --input rerun/results.json
pdflatex -interaction=nonstopmode -halt-on-error article.tex
pdflatex -interaction=nonstopmode -halt-on-error article.tex
pdflatex -interaction=nonstopmode -halt-on-error article.tex
\end{verbatim}
The full verification writes \texttt{results.json}, three CSV files, and
prints its exact-check count and diagnostics. The table generator turns
those retained outputs into the three small TeX table fragments input by
the article. A quicker run uses \texttt{--quick} and stops the coefficient
diagnostics at $n=4096$; it should not be used to regenerate the full tables
without restoring the full data afterward.

The supplied PDF was built from the supplied TeX and table fragments.
The exact algebra checks were run; the numerical diagnostics were run;
the article's analytic proofs were not checked by a proof assistant.

\clearpage
% ed. (2026-09-29): widest label 99, since four editorial entries follow the delivered ones.
\begin{thebibliography}{99}
\bibitem{repo-critical}
V. Reshetnikov, ProveIt repository, research manuscript
\emph{Beyond Finite-Action Folds: Critical Hahn Transseries, Stable Sector
Asymptotics, and Sharp Action Budgets}, especially research Question~1.
Snapshot \texttt{ebd8344bca77d8352cd745b7df374618a290a029}, inspected
September 29, 2026.
\href{https://github.com/VladimirReshetnikov/ProveIt/blob/ebd8344bca77d8352cd745b7df374618a290a029/Analysis/Transseries/docs/series-and-transseries/Critical_Hahn_Transseries_Beyond_Finite_Action_Folds/article.tex}{Pinned manuscript source}.

\bibitem{repo-index}
V. Reshetnikov, ProveIt, \emph{Series and transseries} documentation index,
same snapshot. The index records provenance, scope, overlap, and
formalization boundaries of the recent research packages.
\href{https://github.com/VladimirReshetnikov/ProveIt/blob/ebd8344bca77d8352cd745b7df374618a290a029/Analysis/Transseries/docs/series-and-transseries/README.md}{Pinned documentation index}.

\bibitem{bgt}
N. H. Bingham, C. M. Goldie, and J. L. Teugels,
\emph{Regular Variation}, Encyclopedia of Mathematics and its
Applications 27, Cambridge University Press, 1987.
See Chapter~1 for uniform convergence, Potter bounds, asymptotic inversion,
and conjugacy; Appendix~5 for calculating de Bruijn conjugates.
\url{https://doi.org/10.1017/CBO9780511721434}.

\bibitem{janson}
S. Janson, \emph{Simply generated trees, conditioned Galton--Watson trees,
random allocations and condensation}, Probability Surveys \textbf{9}
(2012), 103--252. The arXiv version has different section numbering from
the published version.
\url{https://doi.org/10.1214/11-PS188};
\url{https://arxiv.org/abs/1112.0510}.

\bibitem{dlmf-polylog}
NIST Digital Library of Mathematical Functions,
\S25.12, especially Eq.~25.12.12, \emph{Polylogarithms}.
\url{https://dlmf.nist.gov/25.12.E12}.
Accessed September 29, 2026.

\bibitem{dlmf-lambert}
NIST Digital Library of Mathematical Functions,
\S4.13, \emph{Lambert $W$-Function}.
\url{https://dlmf.nist.gov/4.13}.
Accessed September 29, 2026.

\bibitem{edgar}
G. A. Edgar, \emph{Transseries for beginners},
arXiv:0801.4877, 2008.
\url{https://arxiv.org/abs/0801.4877}.
% ed. (2026-09-29): bibliography entries for the editorial notes (packages of this series).
\bibitem{ed:lce}
\emph{The Logarithmic Critical Endpoint}, ProveIt research package filed in
batch 48, directory \path{Logarithmic_Critical_Endpoint_Lambert_Charts/}
beside the manuscript of \cite{repo-critical}. Editorial addition.

\bibitem{ed:mct}
\emph{Marginal Critical Transseries}, ProveIt research package filed in
batch 48, directory \path{Marginal_Critical_Transseries_Action_Budgets/}
beside the manuscript of \cite{repo-critical}. Editorial addition.

\bibitem{ed:cct}
\emph{Confluent Critical Transseries Across the Exponent-Two Boundary},
ProveIt research package filed in batch 49 (after the inspected snapshot),
directory \path{Confluent_Critical_Transseries_Exponent_Two_Boundary/}
beside the manuscript of \cite{repo-critical}. Editorial addition.

\bibitem{ed:sge}
\emph{Through the Stable--Gaussian Endpoint}, ProveIt research package filed
in batch 49 (after the inspected snapshot), directory
\path{Stable_Gaussian_Endpoint_Uniform_Critical_Transseries/} beside the
manuscript of \cite{repo-critical}. Editorial addition.
\end{thebibliography}
\end{document}
